\section{Introduction}
\label{sec:intro}

The cost of a diffusion-transformer forward pass grows superlinearly with
resolution: token count grows quadratically with edge length, and
attention cost grows again in token count. This cost structure makes
the title question worth stating precisely. At high noise, can the
\emph{gradient} of an adapter be computed on a downscaled latent as a
drop-in substitute for the native one, with the same model, the same
objective, and the same $\sigma$ distribution, and with only the
spatial grid changed on some steps?

There is a well-studied reason to expect an affirmative answer. As the
noise level $\sigma$ grows, per-frequency signal-to-noise falls, and
high spatial frequencies drop below the noise floor first, so that
above a spectral crossover a downscaled latent carries \emph{almost}
the full surviving signal. Studies of this coupling, however, are
confined almost entirely to the \emph{inference} side: scale-wise
distillation \citep{starodubcev2025swd}, spectral progressive
diffusion \citep{xiao2026spd}, and spectrally-guided noise schedules
\citep{esteves2026spectral} run high-noise \emph{sampling} steps at
reduced resolution and judge the premise by the generated output.
Where resolution is reduced during training
\citep{jin2024pyramidal,karras2018progressive,chen2024pixartalpha,hoogeboom2023simple},
the low-resolution stage is given its own, modified objective, and
success is again judged from the final samples. By the spectral
premise, substitution is safe \emph{whenever the noise has masked the
frequencies the coarse grid loses}.

None of this evidence resolves the question at \emph{training} time:
whether a downscaled step delivers the native gradient direction
under an unchanged objective. Output-level success observes the
substitution only through the confounds of a full training or
sampling pipeline, and a small difference in a quality score does not
imply a small difference in what a training step \emph{learns}.
Measuring that gap directly is the route this paper takes. What a
training step consumes is the expected adapter gradient, and
\emph{demotion}, our term for the training-time substitution of a
downscaled, re-encoded latent for the native one, perturbs more than
the network's input: the regression target carries the clean image at
unit weight at every noise level, and the compute graph itself
changes with the grid. Measuring the gap honestly turns out to be a
metrology problem. The natural estimator, a redraw floor minus a
single demoted-arm cosine, is biased by finite-draw variance, and the
bias grows as the demoted grid shrinks, manufacturing exactly the
token-count-ordered floor signature under test. Every claim below is
therefore stated in debiased units, at the instrument's measured
resolution.

The contributions:

\begin{enumerate}
\item \textbf{Two accounts, stated in the same units}
(\S\ref{sec:theory}). We define the demotion gap at the gradient level,
then derive in those units (a)~the gradient gap implied by the
existing spectral argument, taken in its strongest
tolerance-parameterized form (SPD, \citealp{xiao2026spd}), which is
\emph{ratio-only} and \emph{floorless} by construction, and (b)~our
own account, which treats demotion as a perturbation of both factors
of $g = J^{\!\top}r$ and reduces, under four named assumptions, to a
route amplitude on a universal mismatch curve \emph{plus} a
$\sigma$-independent graph floor. The two differ structurally before
any measurement (\S\ref{sec:spectral}), and every experiment is a
point of disagreement between them.
\item \textbf{A debiased measurement that scores both, and a
predictive tool} (\S\ref{sec:evidence}; the resulting
$(\text{route}, \sigma)$ map in \S\ref{sec:map}). Measured with
per-arm self-floors, a draw-count extrapolation, and a pre-specified
margin whose certification power is set, though not defined, by the
instrument's resolution, the map reads as follows: the mildest route
($1024{\to}896$) is
indistinguishable from the re-encode control over the interior
high-noise window $\sigma \in (0.5, 0.94]$; the deeper $1024{\to}768$
route retains a narrow low-excess window ($0.65 < \sigma < 0.95$)
atop a persistent floor; and no noise level rescues a $2\times$
downscale. No member of the spectral tolerance family reproduces this
map; in particular, none locates the $1024{\to}768$ window, nor the
absence of any window for $1024{\to}512$. Head to head on the same
curves, the spectral family misses at $0.147$--$0.355$ RMSE, whereas
our account predicts held-out routes at ${\sim}0.07$--$0.09$, better
than an oracle fit on the held-out data itself, and the floor law
correctly predicts both held-out floors.
\item \textbf{Selective low-resolution training} (\S\ref{sec:practical}).
Gating substitution on the measured-safe window gives an opt-in
$\sigma$-conditional trainer realizing $-15.1\%$ forward FLOPs and
$-14.6\%$ wall-clock at a fixed step budget, with a weight-space
endpoint footprint a factor of $2$--$3$ inside the training-seed
lottery, rising to endpoint cosine $0.75$ when demotion is
late-scheduled (Table~\ref{tab:yardstick}). Generation-quality
non-inferiority at production scale remains the stated open gate.
\end{enumerate}

\section{Related work}
\label{sec:related}

\paragraph{Flow matching and DiT.}
Flow matching \citep{lipman2023flow} regresses a velocity field along
prescribed probability paths between noise and data; rectified flow
\citep{liu2023rectified} takes the paths to be straight lines, giving
the linear noising $z_\sigma = (1-\sigma)x + \sigma\epsilon$ and the
constant target $v = \epsilon - x$ used throughout, the training
objective of current large text-to-image systems
\citep{esser2024scaling}. The Diffusion Transformer
\citep{peebles2023dit} replaces the U-Net denoiser with a transformer
over patchified latent tokens, so a spatial grid enters the network
only as a token sequence: token count grows quadratically with edge
length and attention cost quadratically again in token count, the cost
structure that makes resolution the expensive axis. Spatial position is
now widely carried by a grid-dependent coordinate system, most commonly
rotary embeddings \citep{su2024roformer}.

\paragraph{Scale-wise and progressive-resolution diffusion.}
Progressive growing \citep{karras2018progressive} and resolution
curricula \citep{chen2024pixartalpha,hoogeboom2023simple} train at
increasing resolution as a schedule, each low-resolution stage carrying
its own stage-specific objective. A more recent family ties resolution
to the noise level itself, on the spectral premise above, and has so
far been studied on the \emph{inference} side. Scale-wise distillation
(SwD) \citep{starodubcev2025swd} states the claim directly (above the
crossover, the noised downscaled latent is close in distribution to the
downscaled noised latent) and samples early steps at reduced
resolution. Spectral Progressive Diffusion (SPD) \citep{xiao2026spd}
gives the argument its sharpest available form: a
tolerance-parameterized per-frequency activation time under a diagonal
Gaussian spectral model, a criterion on the per-band output of the
\emph{Bayes-optimal predictor} rather than merely on input SNR, from
which an inference-time resolution schedule follows. Spectrally-guided
noise schedules \citep{esteves2026spectral} apply the same reasoning
to schedule design; where the premise does enter training, as in
pyramidal flow matching \citep{jin2024pyramidal}, the objective is
restaged per pyramid level rather than left native.

\paragraph{Positional extension.}
Position interpolation \citep{chen2023pi}, banded/frequency-selective
variants (YaRN, \citealp{peng2024yarn}), and diffusion-specific dynamic
position extrapolation \citep{issachar2025dype,zhao2026sigma} rescale
rotary coordinates to evaluate a model at unseen grid sizes. We use
exact positional interpolation as a \emph{causal intervention} that
matches the downscaled grid's relative phase geometry to native, in
order to decompose the resolution-sensitivity floor, and adopt the
$\sigma$-gated band schedule of \citet{zhao2026sigma} as a
training-time refinement.
